\SourcePage{58}
\SGAsection{2}{可预表的余纤维群胚}
设 $\mathcal A$ 是 $\mathcal C_\Lambda$ 上的一个\Term{cofibered-groupoid}{余纤维群胚}，并有一个\Term{minimally-versal}{极小万有}\Term{formal-object}{形式对象} $\xi$。\Term{generic}{一般}说来，单凭 $\xi$ 不能恢复 $\mathcal A$（在 $\mathcal C_\Lambda$-\Term{equivalence}{等价}意义下）。还须知道，两个映射 $f,g:R\to S$ 何时给出一个 $\mathcal C_\Lambda$-\Term{isomorphism}{同构} $\xi(f)\to\xi(g)$，其中 $R=P(\xi)$。作为 1.11 和 1.12 的应用，现在来考察在 $\mathcal C_\Lambda$-\Term{equivalence}{等价}意义下“恢复 $\mathcal A$”的问题。
\SourcePage{59}
事实上，我们将得到一个相当满意的答案，适用于很广的一类\Term{cofibered-groupoid}{余纤维群胚}。
实际中常见的许多\Term{cofibered-groupoid}{余纤维群胚}具有一种比\Term{semi-homogeneous}{半齐性}略强的性质，其明确定义见 2.5。为了研究它们，宜用“\Term{smooth-functor}{光滑函子}”的语言重新表述前面的工作。本节始终沿用上一节的所有记号和术语。
\begin{SGAresult}{定义}{2.1}
设 $\varphi:\mathcal A\to\mathcal B$ 是 $\mathcal C_\Lambda$ 上\Term{cofibered-groupoid}{余纤维群胚}的\Term{functor}{函子}。若 $\mathcal B$ 中的每一个\Term{surjectivity}{满射} $\beta:b'\to\varphi(a)$ 都能补成\Term{commutative-diagram}{交换图}
\[
\begin{tikzcd}
\varphi(a')\arrow[r]\arrow[dr,"\varphi(\alpha)"']&b'\arrow[d,"\beta"]\\
&\varphi(a)
\end{tikzcd}\tag{*}
\]
其中 $\alpha:a'\to a$ 是 $\mathcal A$ 中的某个\Term{arrow}{箭头}，则称 $\varphi$ 为\Term{smooth-functor}{光滑的}。若它\Term{smooth-functor}{光滑}，且
$\varphi(k[\epsilon]):[\mathcal A(k[\epsilon])]\to[\mathcal B(k[\epsilon])]$ 是\Term{bijection}{双射}，则称它为\Term{minimally-smooth}{极小光滑的}。
\end{SGAresult}
\textbf{注记 2.2(a)。} 设 $\varphi:\mathcal A\to\mathcal B$ 是 $\mathcal C_\Lambda$ 上\Term{cofibered-groupoid}{余纤维群胚}的\Term{functor}{函子}。注意下列条件\Term{equivalence}{等价}：“$\varphi$ \Term{smooth-functor}{光滑}”；“$\mathcal B$ 中每个\Term{surjectivity}{满射} $\beta:b'\to\varphi(a)$ 都能补成\Term{commutative-diagram}{交换图} $(*)$，且 $\varphi(a')\to b'$ 是 $\mathcal C_\Lambda$-\Term{isomorphism}{同构}”；“$\mathcal B$ 中每个\Term{small-prolongation}{小延拓} $\beta:b'\to\varphi(a)$ 都能补成\Term{commutative-diagram}{交换图} $(*)$”。

\textbf{注记 2.2(b)。} 对 $\mathcal C_\Lambda$ 上任一\Term{cofibered-groupoid}{余纤维群胚} $\mathcal A$，自然\Term{functor}{函子} $\mathcal A\to[\mathcal A]$ 是\Term{minimally-smooth}{极小光滑的}。

\textbf{注记 2.2(c)。} 对 $\widehat{\mathcal C}_\Lambda$ 的每个\Term{object}{对象} $R$，以 $h_R:\mathcal C_\Lambda\to(\mathrm{Ens})$ 表示由
$h_R(S)=\operatorname{Hom}_{\Lambda\text{-alg}}(R,S)$ 定义的\Term{functor}{函子}。它给出 $\mathcal C_\Lambda$ 上一个\Term{cofibered-groupoid}{余纤维离散群胚} $\underline h_R$，而 $\widehat{\mathcal C}_\Lambda$ 中的映射 $R\to S$
\SourcePage{60}
诱导\Term{cofibered-groupoid}{余纤维群胚}的\Term{functor}{函子} $\underline h_S\to\underline h_R$。由 $\widehat{\mathcal C}_\Lambda$ 中映射 $R\to S$ 诱导的\Term{functor}{函子} $\underline h_S\to\underline h_R$ \Term{smooth-functor}{光滑}，当且仅当 $S$ 是 $R$ 上的\Term{formal-power-series-ring}{形式幂级数环}。

\textbf{注记 2.2(d)。} 给定 $R\in\operatorname{ob}(\widehat{\mathcal C}_\Lambda)$，以 $(R,\mathcal C_\Lambda)$ 表示 $\widehat{\mathcal C}_\Lambda$ 中以 $R$ 为源、以 $\mathcal C_\Lambda$ 的\Term{object}{对象}为靶的\Term{arrow-category}{箭头范畴}。设 $\mathcal A$ 是 $\mathcal C_\Lambda$ 上的\Term{cofibered-groupoid}{余纤维群胚}。给定 $\xi\in\operatorname{ob}\widehat{\mathcal A}$，$P(\xi)=R$，设 $\mathcal A'\to(R,\mathcal C_\Lambda)$ 为 $\mathcal A\to\mathcal C_\Lambda$ 沿\Term{canonical}{典范}\Term{functor}{函子} $(R,\mathcal C_\Lambda)\to\mathcal C_\Lambda$ 的拉回，选择一个\Term{section}{截面} $\xi^\#$，使 $\xi^\#(1_R)=\xi$。（这只是说，对每个映射 $f:R\to S$ 选择一个 $f$-映射 $\xi\to\xi(f)$。）于是由 $(R\xrightarrow f S)\mapsto\xi(f)$ 得到 $\mathcal C_\Lambda$ 上的\Term{functor}{函子} $\xi^\#:\underline h_R\to\mathcal A$。当然，在\Term{functor}{函子}的自然\Term{equivalence}{等价}意义下，它不依赖于 $\xi^\#$ 的选择。“$\xi^\#:\underline h_R\to\mathcal A$（极小）\Term{smooth-functor}{光滑}”意味着 $\mathcal A$ 中每一个\Term{surjectivity}{满射} $\alpha:a'\to\xi(f)$ 都能补成\Term{commutative-diagram}{交换图}
\[
\begin{tikzcd}
\xi(f')\arrow[r]\arrow[dr,"\xi(\ )"']&a'\arrow[d,"\alpha"]\\
&\xi(f)
\end{tikzcd}
\]
（且 $h_R(k[\epsilon])\to[\mathcal A(k[\epsilon])]$ 是\Term{bijection}{双射}），这恰好意味着 $\xi$ 是 $\mathcal A$ 在 $\mathcal C_\Lambda$ 上的（\Term{minimally-versal}{极小}）\Term{versal}{万有}\Term{formal-object}{形式对象}。所以，$\mathcal A$ 有（\Term{minimally-versal}{极小}）\Term{versal}{万有}\Term{formal-object}{形式对象}，\Term{equivalence}{等价}于存在某个 $R\in\operatorname{ob}(\widehat{\mathcal C}_\Lambda)$ 以及 $\mathcal C_\Lambda$ 上一个（\Term{minimally-smooth}{极小}）\Term{smooth-functor}{光滑函子} $\underline h_R\to\mathcal A$。

以下关于\Term{smooth-functor}{光滑性}的\Term{generic}{一般}性质很容易验证。
\begin{SGAresult}{命题}{2.3}
设 $\mathcal A\xrightarrow\varphi\mathcal B\xrightarrow\psi\mathcal C$ 是 $\mathcal C_\Lambda$ 上\Term{cofibered-groupoid}{余纤维群胚}的\Term{functor}{函子}。
(a) 若 $\varphi,\psi$ 都\Term{smooth-functor}{光滑}，则 $\psi\circ\varphi$ 也\Term{smooth-functor}{光滑}。
(b) 若 $\psi\circ\varphi$ \Term{smooth-functor}{光滑}，且 $\varphi$ 在\Term{object}{对象}的\Term{isomorphism-class}{同构类}上\Term{surjectivity}{满射}，
\SourcePage{61}
则 $\psi$ \Term{smooth-functor}{光滑}。
\end{SGAresult}
\textbf{注记 2.4。} 设 $\mathcal A$ 是 $\mathcal C_\Lambda$ 上的\Term{cofibered-groupoid}{余纤维群胚}，有\Term{versal}{万有}\Term{formal-object}{形式对象} $\xi$。由 2.2(d)，这意味着 $\widetilde\xi:\underline h_R\to\mathcal A$ \Term{smooth-functor}{光滑}，其中 $R=P(\xi)$。因此由 2.3(b)，下列条件\Term{equivalence}{等价}：“$\mathcal A$ 在 $\mathcal C_\Lambda$ 上（\Term{formally-smooth}{形式光滑}）（参见 1.13），即 $R$ 是 $\Lambda$ 上的\Term{formal-power-series-ring}{形式幂级数代数}”；“复合\Term{functor}{函子} $\underline h_R\xrightarrow{\widetilde\xi}\mathcal A\xrightarrow P\mathcal C_\Lambda\ (=\underline h_\Lambda)$ \Term{smooth-functor}{光滑}”；“$\mathcal A\xrightarrow P\mathcal C_\Lambda$ 是\Term{smooth-functor}{光滑函子}”。

现在考察一种比\Term{semi-homogeneous}{半齐性}更强的性质。
\begin{SGAresult}{定义}{2.5}
设 $\mathcal A$ 是 $\mathcal C_\Lambda$ 上的\Term{cofibered-groupoid}{余纤维群胚}。若每个图
\[
\begin{tikzcd}&a'\arrow[d]\\b'\arrow[r]&a\end{tikzcd}
\]
其中 $a'\to a$ 为\Term{surjectivity}{满射}，都在 $\mathcal A$ 中有\Term{fiber-product}{纤维积} $a'\times_a b'$，并且
\[P(a'\times_a b')=P(a')\times_{P(a)}P(b'),\]
则称 $\mathcal A$ 为\Term{homogeneous}{齐性的}。
\end{SGAresult}
\textbf{注记 2.6(a)。} 设 $\mathcal A$ 是 $\mathcal C_\Lambda$ 上的\Term{cofibered-groupoid}{余纤维群胚}，考虑 $\mathcal A$ 中的图 $(*)$ 及其在 $\mathcal C_\Lambda$ 中的像
\[
\begin{tikzcd}&a'\arrow[dl]\arrow[dr]\\a_1\arrow[dr]&&a_2\arrow[dl]\\&a\end{tikzcd}
\qquad
\begin{tikzcd}&R_1\times_R R_2\arrow[dl,"\pi_1"']\arrow[dr,"\pi_2"]\\R_1\arrow[dr]&&R_2\arrow[dl]\\&R.\end{tikzcd}
\]
这里 $\pi_i$ 为\Term{projection}{投影}。若\Term{fiber-product}{纤维积} $a_1\times_a a_2$ 在 $\mathcal A$ 中存在，则得到唯一的、与\Term{projection}{投影}相容的映射 $\varphi:a'\to a_1\times_a a_2$。若 $P(a_1\times_a a_2)=R_1\times_R R_2$，则 $P(\varphi)$ 是 $R_1\times_R R_2$ 的\Term{identity-map}{恒等映射}，因而 $\varphi$ 是 $\mathcal C_\Lambda$-\Term{isomorphism}{同构}，即 $a'\simeq a_1\times_a a_2$。换言之，若 $\mathcal A$ 为 $\mathcal C_\Lambda$ 上的\Term{homogeneous}{齐性}\Term{cofibered-groupoid}{余纤维群胚}，则只要 $R_1\to R$ 在 $\mathcal C_\Lambda$ 中\Term{surjectivity}{满射}，每个图 $(*)$ 都定义一个\Term{fiber-product}{纤维积} $a_1\times_a a_2$。这立即说明
\SourcePage{62}
\Term{homogeneous}{齐性}\Term{cofibered-groupoid}{余纤维群胚}当然也是\Term{semi-homogeneous}{半齐性的}。

\textbf{注记 2.6(b)。} 考虑\Term{category}{范畴}与\Term{functor}{函子}构成的图 $X\xrightarrow\varphi Z\xleftarrow\psi Y$。通常的\Term{fiber-product}{纤维积} $X\times_ZY$ 是如下\Term{category}{范畴}：\Term{object}{对象}是满足 $\varphi(x)=\psi(y)$ 的一对 $(x,y)$；映射 $(x,y)\to(x',y')$ 是满足 $\varphi(\alpha)=\psi(\beta)$ 的一对 $\alpha:x\to x'$、$\beta:y\to y'$。$2$-\Term{category}{范畴}意义下的\Term{fiber-product}{纤维积} $X\mathbin{\overset{2}{\times}}_ZY$ 的\Term{object}{对象}是三元组 $(x,y;\gamma)$，其中 $\gamma:\varphi(x)\xrightarrow\sim\psi(y)$ 为\Term{isomorphism}{同构}；映射 $(x,y;\gamma)\to(x',y';\gamma')$ 是使下图交换的一对 $x\xrightarrow\alpha x'$、$y\xrightarrow\beta y'$：
\[
\begin{tikzcd}
\varphi(x)\arrow[r,"\varphi(\alpha)"]\arrow[d,"\gamma"']&\varphi(x')\arrow[d,"\gamma'"]\\
\psi(y)\arrow[r,"\psi(\beta)"']&\psi(y').
\end{tikzcd}
\]
可以用\Term{category}{范畴}的\Term{fiber-product}{纤维积}（在 $2$-\Term{category}{范畴}意义下）重新表述（\Term{semi-homogeneous}{半}）\Term{homogeneous}{齐性}\Term{cofibered-groupoid}{余纤维群胚}及\Term{smooth-functor}{光滑函子}的概念。设 $P:\mathcal A\to\mathcal C_\Lambda$ 为\Term{cofibered-category}{余纤维范畴}，考虑 $\mathcal C_\Lambda$ 中的\Term{object}{对象} $R_1\times_RR_2$。对\Term{projection}{投影} $\pi_i:R_1\times_RR_2\to R_i$（$i=1,2$），若选择\Term{direct-image}{直接像}\Term{functor}{函子} $(\pi_i)_*:\mathcal A(R_1\times_RR_2)\to\mathcal A(R_i)$，便得到\Term{functor}{函子}
\[\mathcal A(R_1\times_RR_2)\longrightarrow\mathcal A(R_1)\mathbin{\overset{2}{\times}}_{\mathcal A(R)}\mathcal A(R_2).\]
由\Term{cofibered-category}{余纤维范畴}的定义可知，在\Term{canonical}{典范}自然\Term{equivalence}{等价}意义下，它不依赖于\Term{direct-image}{直接像}\Term{functor}{函子}的选择。由定义也可知，$\mathcal C_\Lambda$ 上的\Term{cofibered-groupoid}{余纤维群胚} $\mathcal A$ 为\Term{homogeneous}{齐性的}，当且仅当每逢 $R_1\to R$ 在 $\mathcal C_\Lambda$ 中\Term{surjectivity}{满射}，上述\Term{functor}{函子}都是\Term{category-equivalence}{范畴等价}；$\mathcal A$ 为\Term{semi-homogeneous}{半齐性的}，当且仅当
\[ [\mathcal A(R_1\times_RR_2)]\longrightarrow
[\mathcal A(R_1)\mathbin{\overset{2}{\times}}_{\mathcal A(R)}\mathcal A(R_2)]\]
\SourcePage{63}
在 $R_1\to R$ \Term{surjectivity}{满射}时为\Term{surjectivity}{满射}，并在 $R=k$、$R_1=k[\epsilon]$ 时为\Term{bijection}{双射}。再设 $\varphi:\mathcal A\to\mathcal B$ 是 $\mathcal C_\Lambda$ 上\Term{cofibered-groupoid}{余纤维群胚}的\Term{functor}{函子}。对 $\mathcal C_\Lambda$ 中每个映射 $R'\to R$，$\varphi$ 诱导一个在（\Term{canonical}{典范}）自然\Term{equivalence}{等价}意义下唯一确定的\Term{functor}{函子}
\[\mathcal A(R')\longrightarrow\mathcal A(R)\mathbin{\overset{2}{\times}}_{\mathcal B(R)}\mathcal B(R').\]
仍然容易看出，$\varphi$ \Term{smooth-functor}{光滑}当且仅当每逢 $R'\to R$ \Term{surjectivity}{满射}，此\Term{functor}{函子}在\Term{object}{对象}的\Term{isomorphism-class}{同构类}上\Term{surjectivity}{满射}。

\textbf{注记 2.6(c)。} 若 $\mathcal A$ 是 $\mathcal C_\Lambda$ 上的\Term{homogeneous}{齐性}\Term{cofibered-groupoid}{余纤维群胚}，则对任意\Term{finite-dimensional}{有限维} $k$-\Term{vector-space}{向量空间} $V,W$，\Term{functor}{函子}
$\mathcal A(k[V\times W])\to\mathcal A(k[V])\times\mathcal A(k[W])$ 为\Term{category-equivalence}{范畴等价}（回忆约定 $\mathcal A(k)=\{e\}$）。特别地，
\[\operatorname{Aut}_{\mathcal A(k[V\times W])}(1_{V\times W})\simeq
\operatorname{Aut}_{\mathcal A(k[V])}(1_V)\times\operatorname{Aut}_{\mathcal A(k[W])}(1_W),\]
其中 $1_V$ 为 $e$ 沿 $\mathcal C_\Lambda$ 中唯一映射 $k\to k[V]$ 的\Term{direct-image}{直接像}。所以，若 $\mathcal A$ 在 $\mathcal C_\Lambda$ 上为\Term{homogeneous}{齐性的}，则 $\operatorname{Aut}(1_\epsilon)$ 和 $[\mathcal A(k[\epsilon])]$ 都\Term{canonical}{典范}地具有 $k$-\Term{vector-space}{向量空间}结构。

\textbf{注记 2.6(d)。} 对 $\widehat{\mathcal C}_\Lambda$ 中任一\Term{object}{对象} $R$，$\mathcal C_\Lambda$ 上的\Term{cofibered-groupoid}{余纤维（离散）群胚} $\underline h_R$ 是\Term{homogeneous}{齐性的}。

若 $\mathcal A$ 是\Term{semi-homogeneous}{半齐性的}，则 $[\mathcal A]$ 也是如此。然而，$\mathcal A$ 的\Term{homogeneous}{齐性}\Term{generic}{一般}并不蕴含 $[\mathcal A]$ 的\Term{homogeneous}{齐性}。事实上，有：
\begin{SGAresult}{命题}{2.7}
设 $\mathcal A$ 是 $\mathcal C_\Lambda$ 上的\Term{homogeneous}{齐性}\Term{cofibered-groupoid}{余纤维群胚}。以下陈述\Term{equivalence}{等价}。
(i) $[\mathcal A]$ 为\Term{homogeneous}{齐性的}，即对每个\Term{small-extension}{小扩张} $R'\to R$，映射
\[[\mathcal A(R'\times_RS)]\longrightarrow[\mathcal A(R')]\times_{[\mathcal A(R)]}[\mathcal A(S)]\]
为\Term{bijection}{双射}。
(ii) 对 $\mathcal A$ 中每个\Term{small-prolongation}{小延拓} $a'\to a$，映射
\[\operatorname{Aut}_{\mathcal A(R')}(a')\longrightarrow\operatorname{Aut}_{\mathcal A(R)}(a)\]
\SourcePage{64}
为\Term{surjectivity}{满射}，其中 $R'=P(a')$、$R=P(a)$。
\end{SGAresult}
\begin{proof}
(i)$\Rightarrow$(ii)。设 $f:R'\to R$ 是 $\mathcal C_\Lambda$ 中的\Term{small-extension}{小扩张}，$a'\to a$ 是 $\mathcal A$ 中的 $f$-映射。给定 $\tau\in\operatorname{Aut}_{\mathcal A(R)}(a)$，考虑 $a'\times_a a'$ 与 $a'\times_a a'_\tau$；其中 $a'_\tau\to a$ 表示复合映射 $a'\to a\xrightarrow\tau a$。由于
\[[\mathcal A(R'\times_RR')]\longrightarrow[\mathcal A(R')]\times_{[\mathcal A(R)]}[\mathcal A(R')]\]
为\Term{bijection}{双射}，故存在\Term{isomorphism}{同构} $u:a'\times_a a'\xrightarrow\sim a'\times_a a'_\tau$；通过向第二因子的\Term{projection}{投影}，$u$ 诱导一个\Term{isomorphism}{同构} $\tau':a'\xrightarrow\sim a'$，它\Term{lifting}{提升} $\tau$。

(ii)$\Rightarrow$(i)。设 $a_i''$（$i=1,2$）为 $\mathcal A(R'\times_RS)$ 中的\Term{object}{对象}，它们在
$[\mathcal A(R')]\times_{[\mathcal A(R)]}[\mathcal A(S)]$ 中的像相同。于是得到如下实线\Term{arrow}{箭头}图，它位于右侧的底图之上：
\[
\begin{tikzcd}[column sep=small]
&a_2''\arrow[dddl]\arrow[dddrr] &\\
&a_1''\arrow[dl]\arrow[dr]\arrow[u,dashed,"\sim"]&\\
a_1'\arrow[ddr]\arrow[d,"\sim"']&&b_1\arrow[ddl]\arrow[dr,"\tau"]\\
a_2'\arrow[ddr]&&&b_2\arrow[ddll]\\
&a_1\arrow[d,"\bar\tau"']&&\\
&a_2&&
\end{tikzcd}
\qquad
\begin{tikzcd}&R'\times_RS\arrow[dl]\arrow[dr]\\R'\arrow[dr]&&S\arrow[dl]\\&R\end{tikzcd}
\]
这里 $\bar\tau:a_1\to a_2$ 是由 $\tau$ 诱导的\Term{isomorphism}{同构}，使右下方块交换。若条件 (ii) 成立，总能更换给定的\Term{isomorphism}{同构} $a_1'\xrightarrow\sim a_2'$，使左侧方块也交换。由 2.6(a)，对 $i=1,2$ 有 $a_i''\simeq a_i'\times_{a_i}b_i$，所以得到\Term{isomorphism}{同构} $a_1''\xrightarrow\sim a_2''$。这证明所考察的映射为\Term{injectivity}{单射}；又因假设 $\mathcal A$ 为\Term{homogeneous}{齐性的}，它也是\Term{surjectivity}{满射}，因而为\Term{bijection}{双射}。
\end{proof}
设 $\mathcal A$ 为 $\mathcal C_\Lambda$ 上的\Term{homogeneous}{齐性}\Term{cofibered-groupoid}{余纤维群胚}，$f:R'\to R$ 为 $\mathcal C_\Lambda$ 中的\Term{small-extension}{小扩张}。有一个\Term{canonical-map}{典范映射}
\SourcePage{65}
\[\widetilde f:R'\times_RR'\longrightarrow k[\operatorname{Ker}f],\qquad
(r_1',r_2')\longmapsto r_1'(0)+(r_1'-r_2'),\]
且 $1\times\widetilde f:R'\times_RR'\to R'\times_k k[\operatorname{Ker}f]$ 是\Term{isomorphism}{同构}。给定 $\mathcal A$ 中两个 $f$-映射 $a_i\to a$（$i=1,2$），则 $a_1\times_a a_2$ 位于 $R'\times_RR'$ 上方，故可考虑
\[\tau(a_1,a_2)=(\widetilde f)_*(a_1\times_a a_2)\in\operatorname{ob}\mathcal A(k[\operatorname{Ker}f]).\]
因为 $1\times\widetilde f$ 是\Term{isomorphism}{同构}，所以有与向第一因子的\Term{projection}{投影}相容的\Term{isomorphism}{同构}
\[a_1\times_a a_2\simeq a_1\times\tau(a_1,a_2).\]
对\Term{homogeneous}{齐性群胚}，1.5(2) 的陈述可写成更明确的\Term{form}{形式}。
\begin{SGAresult}{\Term{lemma}{引理}}{2.8}
设 $\mathcal A$ 是 $\mathcal C_\Lambda$ 上的\Term{homogeneous}{齐性}\Term{cofibered-groupoid}{余纤维群胚}，$\alpha_i:a_i\to a$（$i=1,2$）是 $\mathcal A$ 中的 $f$-映射。则：
(1) 存在 $\mathcal A$ 中的 $\rho:a_1\to a_2$ 满足 $\alpha_2\rho=\alpha_1$，当且仅当 $\mathcal A$ 中存在\Term{arrow}{箭头} $a_1\to\tau(a_1,a_2)$。
(2) 存在满足 $\alpha_2\rho=\alpha_1$ 及 $P(\rho)=\operatorname{id}_{R'}$ 的 $\rho:a_1\to a_2$，当且仅当 $\tau(a_1,a_2)=0$ 于 $[\mathcal A(k[\operatorname{Ker}f])]$ 中。
\end{SGAresult}
\begin{proof}
同时证明两者。存在满足 $\alpha_2\rho=\alpha_1$ 的 $\rho:a_1\to a_2$（且 $P(\rho)=\operatorname{id}_{R'}$），\Term{equivalence}{等价}于存在 $h:a_1\to a_1\times_a a_2$ 满足 $\operatorname{pr}_1h=\operatorname{id}_{a_1}$（且 $P(h)$ 为\Term{diagonal-map}{对角映射}），\Term{equivalence}{等价}于存在 $g:a_1\to a_1\times\tau(a_1,a_2)$ 满足 $\operatorname{pr}_1g=\operatorname{id}_{a_1}$（且 $P(g):R'\to R'\times_k k[\operatorname{Ker}f]$ 为 $r'\mapsto(r',r'(0))$），\Term{equivalence}{等价}于 $\mathcal A$ 中存在\Term{arrow}{箭头} $\gamma:a_1\to\tau(a_1,a_2)$（且 $P(\gamma):R'\to k[\operatorname{Ker}f]$ 为 $r'\mapsto r'(0)$）。而存在这样的 $\gamma$，使 $P(\gamma)$ 通过 $R'\to k\to k[\operatorname{Ker}f]$ \Term{resolution}{分解}，当且仅当 $\tau(a_1,a_2)=0$。
\end{proof}
\begin{SGAresult}{命题}{2.9}
(a) 设 $\mathcal A\xrightarrow\varphi\mathcal B\xrightarrow\psi\mathcal C$ 为 $\mathcal C_\Lambda$ 上\Term{cofibered-groupoid}{余纤维群胚}的\Term{functor}{函子}。假设：
(i) $\mathcal B$ 为\Term{semi-homogeneous}{半齐性的}，$\mathcal C$ 为\Term{homogeneous}{齐性的}；
(ii) $\psi(k[\epsilon]):[\mathcal B(k[\epsilon])]\to[\mathcal C(k[\epsilon])]$ 为\Term{injectivity}{单射}。
则复合\Term{functor}{函子} $\psi\circ\varphi$ \Term{smooth-functor}{光滑}，当且仅当 $\varphi,\psi$ 都\Term{smooth-functor}{光滑}。
\SourcePage{66}
(b) 设 $\varphi:\mathcal A\to\mathcal B$ 是 $\mathcal C_\Lambda$ 上\Term{cofibered-groupoid}{余纤维群胚}的\Term{smooth-functor}{光滑函子}。若 $\mathcal A$ 为\Term{semi-homogeneous}{半齐性的}，$\mathcal B$ 为\Term{homogeneous}{齐性的}，则 $\widehat\varphi:\widehat{\mathcal A}\to\widehat{\mathcal B}$ \Term{smooth-functor}{光滑}；也就是说，对 $\widehat{\mathcal B}$ 中任一\Term{surjectivity}{满射} $\beta:b'\to\varphi(a)$，存在 $\widehat{\mathcal A}$ 中的 $\alpha:a'\to a$ 和一个 $\widehat{\mathcal C}_\Lambda$-\Term{isomorphism}{同构} $\rho:\varphi(a')\xrightarrow\sim b'$，使 $\beta\rho=\varphi(\alpha)$。
\end{SGAresult}
\begin{proof}
(a) 假设 $\psi\circ\varphi$ \Term{smooth-functor}{光滑}。若 $\varphi$ \Term{smooth-functor}{光滑}，则由 2.3(b)，$\psi$ 也\Term{smooth-functor}{光滑}，所以只须证明 $\varphi$ \Term{smooth-functor}{光滑}。给定 $\mathcal B$ 中的\Term{small-prolongation}{小延拓} $\beta:b'\to\varphi(a)$。由于复合\Term{functor}{函子} $\psi\circ\varphi$ \Term{smooth-functor}{光滑}，存在 $\mathcal A$ 中的\Term{arrow}{箭头} $\alpha:a'\to a$，以及一个 $\mathcal C_\Lambda$-\Term{isomorphism}{同构} $\rho:\psi(b')\xrightarrow\sim\psi\varphi(a')$，使图
\[
\begin{tikzcd}
\psi(b')\arrow[rr,"\rho"]\arrow[dr,"\psi(\beta)"']&&\psi\varphi(a')\arrow[dl,"\psi\varphi(\alpha)"]\\
&\psi\varphi(a)&
\end{tikzcd}
\]
交换。由于 $\mathcal B$ 为\Term{semi-homogeneous}{半齐性的}，可选择\Term{commutative-diagram}{交换图}
\[
\begin{tikzcd}&b''\arrow[dl]\arrow[dr]\\b'\arrow[dr,"\beta"']&&\varphi(a')\arrow[dl,"\varphi(\alpha)"]\\&\varphi(a)\end{tikzcd}
\qquad\text{位于}\quad
\begin{tikzcd}&R'\times_RR'\arrow[dl]\arrow[dr]\\R'\arrow[dr]&&R'\arrow[dl]\\&R\end{tikzcd}
\quad\text{之上，}
\]
其中 $R=P(a)$、$R'=P(b')=P(a')$。因为 $\mathcal C$ 为\Term{homogeneous}{齐性的}，由 2.6(a)，
$\psi(b'')\simeq\psi(b')\times_{\psi\varphi(a)}\psi\varphi(a')$；由 2.8(2)，
\[0=\tau(\psi\varphi(a'),\psi(b'))=(\widetilde f)_*\psi(b''),\]
故由假设 (ii)，$(\widetilde f)_*(b'')=0$。由 1.5(2)，存在映射 $\gamma:b'\to\varphi(a')$ 使 $\beta=\varphi(\alpha)\gamma$。

(b) 给定 $\widehat{\mathcal B}$ 中\Term{surjectivity}{满射} $\beta:b'\to\varphi(a)$。这意味着给定了 $\mathcal B$ 中一列\Term{commutative-diagram}{交换图}
\SourcePage{67}
\[
\begin{tikzcd}
\cdots\arrow[r]&b_3'\arrow[r,"\beta_2'"]\arrow[d,"\beta_3"']&b_2'\arrow[r,"\beta_1'"]\arrow[d,"\beta_2"']&b_1'\arrow[d,"\beta_1"]\\
\cdots\arrow[r]&\varphi(a_3)\arrow[r,"\varphi(\alpha_2)"']&\varphi(a_2)\arrow[r,"\varphi(\alpha_1)"']&\varphi(a_1),
\end{tikzcd}
\]
其中对所有 $\nu$，$\alpha_\nu,\beta_\nu,\beta_\nu'$ 均为\Term{surjectivity}{满射}，且 $\beta=\varprojlim\beta_\nu$。为证明 $\widehat\varphi$ \Term{smooth-functor}{光滑}，只须证明如下实线\Term{arrow}{箭头}\Term{commutative-diagram}{交换图}（其中 $\alpha_n,\beta_n,\beta_{n+1},\beta_n'$ \Term{surjectivity}{满射}）能添上虚线\Term{arrow}{箭头}而仍交换：
\[
\begin{tikzcd}[column sep=small]
\varphi(a_{n+1}')\arrow[rrr,dashed,"\varphi(\ )"]\arrow[dr,dashed]\arrow[ddr,dashed,"\varphi(\ )"']&&&\varphi(a_n')\arrow[dl,"\sim"']\arrow[ddl,"\varphi(\alpha_n')"]\\
&b_{n+1}'\arrow[r,"\beta_n'"]\arrow[d,"\beta_{n+1}"']&b_n'\arrow[d,"\beta_n"]&\\
&\varphi(a_{n+1})\arrow[r,"\varphi(\alpha_n)"']&\varphi(a_n)&
\end{tikzcd}
\]
按假设，$\mathcal A$ 为\Term{semi-homogeneous}{半齐性的}，且 $\alpha_n':a_n'\to a_n$ \Term{surjectivity}{满射}，所以存在 $\mathcal A$ 中\Term{commutative-diagram}{交换图}，位于相应的底图之上：
\[
\begin{tikzcd}&a'\arrow[dl,"\gamma"']\arrow[dr,"\gamma'"]\\a_{n+1}\arrow[dr,"\alpha_n"']&&a_n'\arrow[dl,"\alpha_n'"]\\&a_n\end{tikzcd}
\qquad
\begin{tikzcd}&R_{n+1}\times_{R_n}R_n'\arrow[dl]\arrow[dr]\\R_{n+1}\arrow[dr]&&R_n'\arrow[dl]\\&R_n.\end{tikzcd}
\]
由 $\mathcal B$ 的\Term{homogeneous}{齐性}，下图为 $\mathcal B$ 中的\Term{fiber-product}{纤维积}图：
\[
\begin{tikzcd}&\varphi(a')\arrow[dl,"\varphi(\gamma)"']\arrow[dr,"\varphi(\gamma')"]\\\varphi(a_{n+1})\arrow[dr,"\varphi(\alpha_n)"']&&\varphi(a_n')\arrow[dl,"\varphi(\alpha_n')"]\\&\varphi(a_n).\end{tikzcd}
\]
\SourcePage{68}
因此得到唯一的映射 $\delta:b_{n+1}'\to\varphi(a')$，满足所需的交换关系。因为 $\beta_{n+1},\beta_n'$ 都\Term{surjectivity}{满射}，$\delta$ 也\Term{surjectivity}{满射}。由 $\varphi$ 的\Term{smooth-functor}{光滑性}，存在 $\mathcal A$ 中的 $\alpha':a_{n+1}'\to a'$ 及 $\mathcal C_\Lambda$-\Term{isomorphism}{同构} $\rho:\varphi(a_{n+1}')\xrightarrow\sim b_{n+1}'$，使 $\delta\rho=\varphi(\alpha')$。于是得到\Term{commutative-diagram}{交换图}
\[
\begin{tikzcd}[column sep=small]
\varphi(a_{n+1}')\arrow[rrr,"\varphi(\gamma'\alpha')"]\arrow[dr,"\rho"]\arrow[ddr,"\varphi(\gamma\alpha')"']&&&\varphi(a_n')\arrow[dl,"\sim"']\arrow[ddl,"\varphi(\alpha_n')"]\\
&b_{n+1}'\arrow[r]\arrow[d]&b_n'\arrow[d]&\\
&\varphi(a_{n+1})\arrow[r]&\varphi(a_n).&
\end{tikzcd}
\]
\end{proof}
\begin{SGAresult}{推论}{2.10}
(a) 设 $\varphi:\mathcal A\to\mathcal B$ 是 $\mathcal C_\Lambda$ 上\Term{cofibered-groupoid}{余纤维群胚}的\Term{smooth-functor}{光滑函子}，$\xi$ 为 $\mathcal A$ 的\Term{versal}{万有}\Term{formal-object}{形式对象}，$\eta$ 为 $\mathcal B$ 的\Term{minimally-versal}{极小万有}\Term{formal-object}{形式对象}。若 $\mathcal B$ 为\Term{homogeneous}{齐性的}，则对 $\widehat{\mathcal B}$ 中任意映射 $\beta:\eta\to\varphi(\xi)$，$P(\xi)$ 经由 $P(\beta):P(\eta)\to P(\varphi(\xi))=P(\xi)$ 成为 $P(\eta)$ 上的\Term{formal-power-series-ring}{形式幂级数环}。
(b) 设 $\mathcal A$ 是 $\mathcal C_\Lambda$ 上的\Term{homogeneous}{齐性}\Term{cofibered-groupoid}{余纤维群胚}，$\xi,\eta$ 是 $\mathcal A$ 的\Term{versal}{万有}\Term{formal-object}{形式对象}。若 $\eta$ 极小，则对 $\widehat{\mathcal A}$ 中任意映射 $\alpha:\eta\to\xi$，$P(\xi)$ 经由 $P(\alpha):P(\eta)\to P(\xi)$ 成为 $P(\eta)$ 上的\Term{formal-power-series-ring}{形式幂级数环}。
\end{SGAresult}
\begin{proof}
(a) 映射 $\beta:\eta\to\varphi(\xi)$ 诱导\Term{commutative-diagram}{交换图}
\[
\begin{tikzcd}
\underline h_{P(\xi)}\arrow[r,"\widetilde\beta"]\arrow[d,"\widetilde\xi"']&\underline h_{P(\eta)}\arrow[d,"\widetilde\eta"]\\
\mathcal A\arrow[r,"\varphi"']&\mathcal B.
\end{tikzcd}
\]
由于 $\widetilde\xi,\varphi$ 都\Term{smooth-functor}{光滑}，$\widetilde\eta\widetilde\beta=\varphi\widetilde\xi$ \Term{smooth-functor}{光滑}。另一方面，
\SourcePage{69}
$\widetilde\eta(k[\epsilon]):h_{P(\eta)}(k[\epsilon])\to[\mathcal B(k[\epsilon])]$ 为\Term{isomorphism}{同构}，所以由 2.9(a)，$\widetilde\beta$ \Term{smooth-functor}{光滑}，这意味着 $P(\xi)$ 为 $P(\eta)$ 上的\Term{formal-power-series-ring}{形式幂级数环}。
(b) 在 (a) 中令 $\mathcal A=\mathcal B$ 即得。
\end{proof}
最后，考虑\Term{cofibered-groupoid}{余纤维群胚}的“\Term{representability}{可表性}”问题。设 $\mathcal C$ 为一个固定\Term{category}{范畴}。$\mathcal C$ 中每个\Term{object}{对象} $R$ 都定义一个\Term{functor}{函子} $h_R:\mathcal C\to(\mathrm{Ens})$，$h_R(S)=\operatorname{Hom}_{\mathcal C}(R,S)$，从而得到一个\Term{fully-faithful}{全忠实}\Term{functor}{函子} $h:\mathcal C^\circ\to\operatorname{Hom}(\mathcal C,(\mathrm{Ens}))$。所谓“$\mathcal C$ 中的\Term{cocategory}{余范畴}（cocategory）”，是指一个图
\[
 h_{R_1}\times_{h_{R_0}}h_{R_1}\xrightarrow c h_{R_1}
 \underset{s_2}{\overset{s_1}{\rightrightarrows}}h_{R_0},
 \qquad \epsilon:h_{R_0}\longrightarrow h_{R_1},
 \qquad s_i\epsilon=\operatorname{id}\quad(i=1,2),
\]
其中 $h_{R_1}\times_{h_{R_0}}h_{R_1}=(h_{R_1})_{s_1}\times_{h_{R_0},s_2}(h_{R_1})$，且满足：
(i) $c$ 有\Term{associative-law}{结合律}，并且对 $i=1,2$ 与 $s_i$ 相容；
(ii) 以下两个复合都为\Term{identity-map}{恒等映射}：
\[
\begin{aligned}
h_{R_1}&\xrightarrow\sim h_{R_0}\times_{h_{R_0}}h_{R_1}
 \xrightarrow{\epsilon\times1}h_{R_1}\times_{h_{R_0}}h_{R_1}\xrightarrow c h_{R_1},\\
h_{R_1}&\xrightarrow\sim h_{R_1}\times_{h_{R_0}}h_{R_0}
 \xrightarrow{1\times\epsilon}h_{R_1}\times_{h_{R_0}}h_{R_1}\xrightarrow c h_{R_1}.
\end{aligned}
\]
若 $(R_0,R_1,s_1,s_2,\epsilon,c)$ 是 $\mathcal C$ 中的\Term{cocategory}{余范畴}，则对 $\mathcal C$ 中\Term{object}{对象} $S$，可定义\Term{category}{范畴} $h_{(R_0,R_1,\ldots)}(S)$，令
\[
\operatorname{ob}h_{(R_0,R_1,\ldots)}(S)=\operatorname{Hom}_{\mathcal C}(R_0,S),\qquad
\operatorname{Fl}h_{(R_0,R_1,\ldots)}(S)=\operatorname{Hom}_{\mathcal C}(R_1,S),
\]
\Term{arrow}{箭头}的\Term{composition-law}{合成律}由 $c:h_{R_1}\times_{h_{R_0}}h_{R_1}\to h_{R_1}$ 给出。于是
$h_{(R_0,R_1,\ldots)}:\mathcal C\to(\text{范畴})$ 是一个\Term{functor}{函子}，从而定义 $\mathcal C$ 上一个\Term{cofibered-category}{余纤维范畴} $\underline h_{(R_0,R_1,\ldots)}$。例如，$\mathcal C$ 中每个\Term{object}{对象} $R$ 都定义一个\Term{discrete-cocategory}{离散余范畴} $\underline h_R=\underline h_{(R,R,s_1,s_2,\epsilon,c)}$，只须令 $s_1,s_2,\epsilon,c$ 全为\Term{identity-map}{恒等映射}。
$\mathcal C$ 中\Term{cocategory}{余范畴}之间的\Term{homomorphism}{同态} $(R_0,R_1,\ldots)\to(R_0',R_1',\ldots)$
\SourcePage{70}
是一对映射 $h_{R_0}\to h_{R_0'}$、$h_{R_1}\to h_{R_1'}$，使对 $\mathcal C$ 中每个\Term{object}{对象} $S$，诱导的映射
\[h_{(R_0,R_1,\ldots)}(S)\longrightarrow h_{(R_0',R_1',\ldots)}(S)\]
为\Term{functor}{函子}。例如，每个\Term{cocategory}{余范畴} $(R_0,R_1,\ldots)$ 都带有一个\Term{cocategory}{余范畴}\Term{homomorphism}{同态} $\epsilon:\underline h_{R_0}\to\underline h_{(R_0,R_1,\ldots)}$。按定义，$\mathcal C$ 中的\Term{cogroupoid}{余群胚}（cogroupoid）是 $\mathcal C$ 中的一个\Term{cocategory}{余范畴} $(R_0,R_1,\ldots)$，使对 $\mathcal C$ 中每个\Term{object}{对象} $S$，$h_{(R_0,R_1,\ldots)}(S)$ 都为\Term{groupoid}{群胚}。换言之，它是一个\Term{cocategory}{余范畴} $(R_0,R_1,\ldots)$，使 $\underline h_{(R_0,R_1,\ldots)}$ 成为 $\mathcal C$ 上的\Term{cofibered-groupoid}{余纤维群胚}。

实际中，\Term{cocategory}{余范畴}或\Term{cogroupoid}{余群胚}按如下方式出现。设 $P:\mathcal F\to\mathcal C$ 为\Term{cofibered-category}{余纤维范畴}，$\xi$ 为 $\mathcal F$ 中固定的\Term{object}{对象}。令 $R_0=P(\xi)$，以 $(R_0,\mathcal C)$ 表示 $\mathcal C$ 中以 $R_0$ 为源的\Term{arrow-category}{箭头范畴}。设 $P:\mathcal F'\to(R_0,\mathcal C)$ 是 $\mathcal F$ 沿\Term{canonical}{典范}\Term{functor}{函子} $(R_0,\mathcal C)\to\mathcal C$ 的拉回，选择 $\mathcal F'$ 在 $(R_0,\mathcal C)$ 上的一个\Term{cocartesian}{余笛卡儿}\Term{section}{截面} $\xi^\#$，使 $\xi^\#(\operatorname{id}_{R_0})=\xi$。（这只是说，对 $\mathcal C$ 中每个\Term{arrow}{箭头} $f:R_0\to S$，选择一个\Term{cocartesian}{余笛卡儿} $f$-映射 $\xi\to\xi^\#(f)$。）定义\Term{functor}{函子}
\[\operatorname{Hom}_{\xi^\#}:(R_0\amalg R_0,\mathcal C)\longrightarrow(\mathrm{Ens}),\qquad
X=(f_1,f_2)\longmapsto\operatorname{Hom}_{\mathcal F(S)}(\xi^\#(f_1),\xi^\#(f_2)),\]
其中 $S$ 为 $f_1$ 的靶（也就是 $f_2$ 的靶）。在（\Term{canonical}{典范}）自然\Term{equivalence}{等价}意义下，此\Term{functor}{函子}不依赖于 $\xi^\#$ 的选择。以下略去此区别，简记 $\xi^\#$ 为 $\xi$。若此\Term{functor}{函子}由 $(R_0\amalg R_0,\mathcal C)$ 中的\Term{object}{对象}
$X_0=(R_0\underset{s_2}{\overset{s_1}{\rightrightarrows}}R_1)$ 表示，即若 $\mathcal F$ 中存在一个 $R_0$-映射 $\varphi:\xi(s_1)\to\xi(s_2)$，使诱导的映射
\[\widetilde\varphi:\operatorname{Hom}_{(R_0\amalg R_0,\mathcal C)}(X_0,X)\longrightarrow\operatorname{Hom}_\xi(X)\]
对所有 $X\in(R_0\amalg R_0,\mathcal C)$ 都为\Term{bijection}{双射}，则\Term{object}{对象} $X_0$ 自然定义 $\mathcal C$ 中的一个\Term{cocategory}{余范畴} $(R_0,R_1,\ldots)$。事实上，$\mathcal C$ 中三个\Term{arrow}{箭头}
\SourcePage{71}
$u,v,w:R_0\to S$ 定义合成
\[\operatorname{Hom}_{\mathcal F(S)}(\xi(v),\xi(w))\times
\operatorname{Hom}_{\mathcal F(S)}(\xi(u),\xi(v))\longrightarrow
\operatorname{Hom}_{\mathcal F(S)}(\xi(u),\xi(w)),\]
从而给出\Term{composition-law}{合成律} $h_{R_1}(S)\times_{h_{R_0}(S)}h_{R_1}(S)\to h_{R_1}(S)$，它对 $S$ 具有\Term{functor}{函子}性。因此\Term{functor}{函子} $\operatorname{Hom}_\xi$ 定义 $\mathcal C$ 中的\Term{cocategory}{余范畴} $(R_0,R_1,\ldots)$，进而定义\Term{cofibered-category}{余纤维范畴} $\underline h_{(R_0,R_1,\ldots)}$。按定义，
\[\operatorname{ob}\underline h_{(R_0,R_1,\ldots)}=coprod_{S\in\operatorname{ob}(\mathcal C)}\operatorname{Hom}_{\mathcal C}(R_0,S).\]
\Term{arrow}{箭头} $(R_0\xrightarrow fS)\to(R_0\xrightarrow{f'}S')$ 是一对 $h:S\to S'$、$u:R_1\to S'$，满足 $us_1=hf$、$us_2=f'$；结构\Term{functor}{函子} $P:\underline h_{(R_0,R_1,\ldots)}\to\mathcal C$ 由 $(R_0\xrightarrow fS)\mapsto S$ 给出。现在，$(\xi,\varphi)$ 定义 $\mathcal C$ 上的\Term{functor}{函子}
\[(\xi,\varphi):\underline h_{(R_0,R_1,\ldots)}\longrightarrow\mathcal F,\]
它把 $(R_0\xrightarrow fS)$ 送到 $\xi(f)$，把\Term{arrow}{箭头}
$[(R_0\xrightarrow fS)\xrightarrow{(h,u)}(R_0\xrightarrow{f'}S')]$ 送到复合映射
\[\xi(f)\xrightarrow{h_*}\xi(hf)\xrightarrow{u(\varphi)}\xi(f').\]
对 $\mathcal C$ 中每个\Term{object}{对象} $S$，\Term{functor}{函子}
$(\xi,\varphi)(S):h_{(R_0,R_1,\ldots)}(S)\to\mathcal F(S)$ 是\Term{fully-faithful}{全忠实的}，这正来自由\Term{functor}{函子} $\operatorname{Hom}_\xi$ 的\Term{representability}{可表性}所给出的 $(R_0\rightrightarrows R_1)$ 的定义。因此，\Term{functor}{函子} $(\xi,\varphi):\underline h_{(R_0,R_1,\ldots)}\to\mathcal F$ 是 $\mathcal C$-\Term{equivalence}{等价}，当且仅当 $\xi$ 是 $\mathcal F$ 在 $\mathcal C$ 上的\Term{quasi-initial}{拟始对象}，即对 $\mathcal F$ 中每个\Term{object}{对象} $a$，都存在\Term{cocartesian}{余笛卡儿}\Term{arrow}{箭头} $\xi\to a$。也可以考虑\Term{functor}{函子}
\[\operatorname{Isom}_\xi:(R_0\amalg R_0,\mathcal C)\longrightarrow(\mathrm{Ens}),\]
\SourcePage{72}
它由 $X=(f_1,f_2)\mapsto\operatorname{Isom}_{\mathcal F(S)}(\xi(f_1),\xi(f_2))$ 给出，其中 $S$ 是 $f_1$（也就是 $f_2$）的靶。若此\Term{functor}{函子}由 $(R_0\amalg R_0,\mathcal C)$ 中\Term{object}{对象} $(R_0\rightrightarrows R_1)$ 表示，则它定义 $\mathcal C$ 中的一个\Term{cogroupoid}{余群胚} $(R_0,R_1,\ldots)$，以及\Term{functor}{函子} $(\xi,\varphi):\underline h_{(R_0,R_1,\ldots)}\to\mathcal F$；当 $\mathcal F$ 为\Term{cofibered-groupoid}{余纤维群胚}时，这一\Term{functor}{函子}是 $\mathcal C$-\Term{equivalence}{等价}，当且仅当 $\xi$ 是 $\mathcal F$ 在 $\mathcal C$ 上的\Term{quasi-initial}{拟始对象}。
\begin{SGAresult}{定义及命题}{2.11}
$\widehat{\mathcal C}_\Lambda$ 中的\Term{cogroupoid}{余群胚} $(R_0,R_1,\ldots)$ 称为\Term{smooth-cogroupoid}{光滑的}，如果它满足以下任一\Term{equivalence}{等价}条件：
(i) $R_1$ 经由 $s_1$（以及 $s_2$）成为 $R_0$ 上的\Term{formal-power-series-ring}{形式幂级数环}；
(ii) $\epsilon:\underline h_{R_0}\to\underline h_{(R_0,R_1,\ldots)}$ 在 $\mathcal C_\Lambda$ 上\Term{smooth-functor}{光滑}。
$\widehat{\mathcal C}_\Lambda$ 中的\Term{minimally-smooth-cogroupoid}{极小光滑余群胚}，或\Term{normalized-cogroupoid}{规范化光滑余群胚}，是指\Term{smooth-cogroupoid}{光滑余群胚} $(R_0,R_1,\ldots)$，使得 $h_{(R_0,R_1,\ldots)}(k[\epsilon])$ 完\Term{totally-disconnected}{全不连通}，即 $s_1(k[\epsilon])=s_2(k[\epsilon])$。
\end{SGAresult}
\begin{proof}
$\epsilon:\underline h_{R_0}\to\underline h_{(R_0,R_1,\ldots)}$ 在 $\mathcal C_\Lambda$ 上\Term{smooth-functor}{光滑}，\Term{equivalence}{等价}于：给定 $\underline h_{(R_0,R_1,\ldots)}$ 中任一\Term{surjectivity}{满射}
$(S',f')\xrightarrow{(h,u)}(S,f)$，存在 $f$ 到 $S'$ 的\Term{lifting}{提升} $f''$，以及\Term{isomorphism}{同构} $(S',f'')\xrightarrow{(1,u')}(S',f')$；\Term{equivalence}{等价}于：给定 $u\in h_{R_1}(S)$ 以及 $s_1(u)$ 到 $S'$ 的\Term{lifting}{提升} $f'$，存在 $u$ 到 $S'$ 的\Term{lifting}{提升} $u'$，使 $s_1(u')=f'$，这里 $S'\to S$ 是 $\mathcal C_\Lambda$ 中任意\Term{surjectivity}{满射}；这又\Term{equivalence}{等价}于 $s_1:\underline h_{R_1}\to\underline h_{R_0}$ 在 $\mathcal C_\Lambda$ 上\Term{smooth-functor}{光滑}。
\end{proof}
\begin{SGAresult}{命题}{2.12}
设 $(R_0,R_1,\ldots)$ 是 $\widehat{\mathcal C}_\Lambda$ 中的\Term{cogroupoid}{余群胚}。
(a) 若 $(R_0,R_1,\ldots)$ 为\Term{smooth-cogroupoid}{光滑的}，则 $\underline h_{(R_0,R_1,\ldots)}$ 为\Term{homogeneous}{齐性的}。
(b) 若 $s_1(k[\epsilon])=s_2(k[\epsilon])$，则以下陈述\Term{equivalence}{等价}：
(i) $(R_0,R_1,\ldots)$ 为\Term{smooth-cogroupoid}{光滑的}；
(ii) $\underline h_{(R_0,R_1,\ldots)}$ 为\Term{homogeneous}{齐性的}；
(iii) $\underline h_{(R_0,R_1,\ldots)}$ 为\Term{semi-homogeneous}{半齐性的}。
\SourcePage{73}
(c) 若 $\underline h_{(R_0,R_1,\ldots)}$ 为\Term{semi-homogeneous}{半齐性的}，则它\Term{equivalence}{等价}于某个\Term{minimally-smooth-cogroupoid}{极小光滑余群胚} $S_*$ 所给出的 $\underline h_{(S_0,S_1,\ldots)}$。
\end{SGAresult}
\begin{proof}
(a) 考虑 $\underline h_{(R_0,R_1,\ldots)}$ 中的图
\[
\begin{tikzcd}
(S_1,f_1)\arrow[dr,"{(h_1,u_1)}"']&&(S_2,f_2)\arrow[dl,"{(h_2,u_2)}"]\\
&(S,f),&
\end{tikzcd}
\]
其中 $h_2:S_2\to S$ 在 $\mathcal C_\Lambda$ 中\Term{surjectivity}{满射}。按假设，$s_1:\underline h_{R_1}\to\underline h_{R_0}$ \Term{smooth-functor}{光滑}，所以 $u_1^{-1}u_2$ \Term{liftable}{可提升}为某个 $v\in h_{R_1}(S_2)$，满足 $s_1(v)=f_2$。由于 $h_2s_2(v)=h_1f_1$，$f_1\times s_2(v)\in h_{R_0}(S_1\times_SS_2)$ 有定义，得到\Term{commutative-diagram}{交换图}
\[
\begin{tikzcd}
&(S_1\times_SS_2,f_1\times s_2(v))\arrow[dl,"\pi_1"']\arrow[dr,"{(\pi_2,v^{-1})}"]&\\
(S_1,f_1)\arrow[dr,"{(h_1,u_1)}"']&&(S_2,f_2)\arrow[dl,"{(h_2,u_2)}"]\\
&(S,f).&
\end{tikzcd}
\]
断言此图在 $\underline h_{(R_0,R_1,\ldots)}$ 中定义\Term{fiber-product}{纤维积} $(S_1,f_1)\times_{(S,f)}(S_2,f_2)$。事实上，对任意\Term{object}{对象} $(T,g)$，有以下等式；为便于排版，令 $Q=S_1\times_SS_2$：
\[
\begin{aligned}
{}&\operatorname{Hom}((T,g),(Q,f_1\times s_2(v)))\\
&=\{(h_1'\times h_2',w_1\times w_2)\in h_T(Q)\times h_{R_1}(Q)\mid
 s_1(w_1\times w_2)=(h_1'\times h_2')g,\\
&\hspace{43mm}s_2(w_1\times w_2)=f_1\times s_2(v)\}\\
&=\{(h_1'\times h_2',w_1,w_2)\in h_T(Q)\times h_{R_1}(S_1)\times h_{R_1}(S_2)\mid
 h_1w_1=h_2w_2,\\
&\hspace{12mm}s_1(w_1)=h_1'g,\ s_1(w_2)=h_2'g,\ s_2(w_1)=f_1,\ s_2(w_2)=s_2(v)\}\\
&=\{(h_1'\times h_2',w_1,w_2)\in h_T(Q)\times h_{R_1}(S_1)\times h_{R_1}(S_2)\mid
 u_1^{-1}u_2h_2(v^{-1}w_2)=h_1w_1,\\
&\hspace{12mm}s_1(w_1)=h_1'g,\ s_2(w_1)=f_1,
\end{aligned}
\]
\SourcePage{74}
\[
\begin{aligned}
{}&\hspace{12mm}s_1(v^{-1}w_2)=h_2'g,\ s_2(v^{-1}w_2)=f_2\}\\
&\xrightarrow\sim\{(h_1'\times h_2',w_1,w_2')\in h_T(Q)\times h_{R_1}(S_1)\times h_{R_1}(S_2)\mid
 u_2h_2w_2'=u_1h_1w_1,\\
&\hspace{12mm}s_1(w_1)=h_1'g,\ s_2(w_1)=f_1,\ s_1(w_2')=h_2'g,\ s_2(w_2')=f_2\}\\
&=\operatorname{Hom}((T,g),(S_1,f_1))
 \times_{\operatorname{Hom}((T,g),(S,f))}\operatorname{Hom}((T,g),(S_2,f_2)).
\end{aligned}
\]
(b) 只须证明 (iii)$\Rightarrow$(i)。由于 $\underline h_{(R_0,R_1,\ldots)}$ 为\Term{semi-homogeneous}{半齐性的}，并且
\[\dim[\underline h_{(R_0,R_1,\ldots)}(k[\epsilon])]=\dim h_{R_0}(k[\epsilon])<\infty,\]
它有一个\Term{minimally-versal}{极小万有}\Term{formal-object}{形式对象}。另一方面，因为
$h_{R_0}(k[\epsilon])\to[\underline h_{(R_0,R_1,\ldots)}(k[\epsilon])]$ 为\Term{bijection}{双射}，$(R_0,1_{R_0})$ 是 $\widehat{\underline h}_{(R_0,R_1,\ldots)}$ 的一个极小\Term{quasi-initial}{拟始对象}，故由 1.10(ii)，它是 $\underline h_{(R_0,R_1,\ldots)}$ 的\Term{minimally-versal}{极小万有}\Term{formal-object}{形式对象}；这意味着 $\epsilon:\underline h_{R_0}\to\underline h_{(R_0,R_1,\ldots)}$ 为\Term{minimally-smooth}{极小光滑的}。

(c) 设 $(S_0,\xi_0)$ 为 $\underline h_{(R_0,R_1,\ldots)}$ 的\Term{minimally-versal}{极小万有}\Term{formal-object}{形式对象}（1.11）。对 $\underline h_{(R_0,R_1,\ldots)}(T)$ 中任意\Term{object}{对象} $\xi:R_0\to T$，考虑从 $\mathcal C_{T\otimes_\Lambda T}$ 到\Term{set}{集合}\Term{category}{范畴}的\Term{functor}{函子} $\operatorname{Isom}(\xi,\xi)$：它把 $R$ 送到 $\xi$ 沿
\[R_0\xrightarrow\xi T\rightrightarrows T\otimes_\Lambda T\longrightarrow R\]
的两个像之间的\Term{isomorphism}{同构}所成的\Term{set}{集合}。它由 $R_1\otimes_{R_0\widehat\otimes_\Lambda R_0}(T\otimes_\Lambda T)$ 表示。取\Term{limit}{极限}便得，$\operatorname{Isom}(\xi_0,\xi_0)$ 由某个 $S_1$ \Term{prorepresentable}{预表示}，且 $\underline h_{(R_0,R_1,\ldots)}$ \Term{equivalence}{等价}于 $\underline h_{(S_0,S_1,\ldots)}$。由 (b)，$S_*$ 为\Term{minimally-smooth-cogroupoid}{极小光滑的}。
\end{proof}
\begin{SGAresult}{定义}{2.13}
$\mathcal C_\Lambda$ 上的\Term{cofibered-groupoid}{余纤维群胚} $\mathcal A$ 称为（极小）\Term{prorepresentable}{可预表的}，若它与 $\widehat{\mathcal C}_\Lambda$ 中某个（\Term{normalized-cogroupoid}{规范化}）\Term{cogroupoid}{余群胚} $(R_0,R_1,\ldots)$ 所给出的 $\underline h_{(R_0,R_1,\ldots)}$ 为 $\mathcal C_\Lambda$-\Term{equivalence}{等价}。
\end{SGAresult}
\begin{SGAresult}{命题}{2.14}
若 $(R_0,R_1,\ldots)$ 为 $\widehat{\mathcal C}_\Lambda$ 中的\Term{normalized-cogroupoid}{规范化余群胚}，则每个诱导 $\mathcal C_\Lambda$-\Term{equivalence}{等价}
$\underline h_{(\alpha_0,\alpha_1)}:\underline h_{(R_0,R_1,\ldots)}\to\underline h_{(R_0,R_1,\ldots)}$
的\Term{homomorphism}{同态} $(\alpha_0,\alpha_1):(R_0,R_1,\ldots)\to(R_0,R_1,\ldots)$ 都是\Term{isomorphism}{同构}。
\SourcePage{75}
特别地，对 $\mathcal C_\Lambda$ 上任意\Term{cofibered-groupoid}{余纤维群胚} $\mathcal A$，\Term{normalized-cogroupoid}{规范化}\Term{prorepresentable}{预表示}若存在，则在\Term{isomorphism}{同构}意义下唯一。
\end{SGAresult}
\begin{proof}
考虑\Term{commutative-diagram}{交换图}
\[
\begin{tikzcd}
\underline h_{(R_0,R_1,\ldots)}\arrow[r,"{\underline h_{(\alpha_0,\alpha_1)}}"]&\underline h_{(R_0,R_1,\ldots)}\\
\underline h_{R_0}\arrow[u,"\epsilon"]\arrow[r,"h_{\alpha_0}"']&\underline h_{R_0}\arrow[u,"\epsilon"'].
\end{tikzcd}
\]
其中 $\underline h_{(\alpha_0,\alpha_1)}$ 为 $\mathcal C_\Lambda$-\Term{equivalence}{等价}。由于 $h_{(R_0,R_1,\ldots)}(k[\epsilon])$ 完\Term{totally-disconnected}{全不连通}，$[\underline h_{(\alpha_0,\alpha_1)}(k[\epsilon])]$ 与 $[\epsilon(k[\epsilon])]$ 均为\Term{bijection}{双射}，故 $h_{\alpha_0}(k[\epsilon])$ 也为\Term{bijection}{双射}。因此 $\alpha_0:R_0\to R_0$ 为满\Term{endomorphism}{自同态}，因而为\Term{automorphism}{自同构}。随后可知 $\alpha_1$ 也为 $R_1$ 的\Term{automorphism}{自同构}。
\end{proof}
\begin{SGAresult}{\Term{lemma}{引理}}{2.15}
设 $\mathcal A,\mathcal B$ 为 $\mathcal C_\Lambda$ 上的\Term{discrete}{离散}\Term{groupoid}{群胚}，$\varphi:\mathcal A\to\mathcal B$ 为\Term{minimally-smooth}{极小光滑}\Term{functor}{函子}。若 $\mathcal A,\mathcal B$ 都为\Term{homogeneous}{齐性的}，则 $\varphi$ 是 $\mathcal C_\Lambda$-\Term{equivalence}{等价}。
\end{SGAresult}
\begin{proof}
须证明对所有 $S\in\operatorname{ob}(\mathcal C_\Lambda)$，$\varphi(S):\mathcal A(S)\to\mathcal B(S)$ 均为\Term{bijection}{双射}。注意，它对所有这样的 $S$ 为\Term{surjectivity}{满射}，对 $S=k[\epsilon]$ 为\Term{bijection}{双射}。对 $S$ 的\Term{length}{长度}归纳。设 $f:S'\to S$ 为 $\mathcal C_\Lambda$ 中的\Term{small-extension}{小扩张}。由于 $\mathcal A,\mathcal B$ 为\Term{homogeneous}{齐性}\Term{discrete}{离散}\Term{groupoid}{群胚}，\Term{canonical-isomorphism}{典范同构}
$1\times\widetilde f:S'\times_SS'\xrightarrow\sim S'\times_k k[\operatorname{Ker}f]$ 诱导\Term{commutative-diagram}{交换图}
\SourcePage{76}
\[
\begin{tikzcd}
\mathcal A(S')\times_{\mathcal A(S)}\mathcal A(S')\arrow[r,"\sim"]\arrow[d]&\mathcal A(S')\times\mathcal A(k[\operatorname{Ker}f])\arrow[d]\\
\mathcal B(S')\times_{\mathcal B(S)}\mathcal B(S')\arrow[r,"\sim"]&\mathcal B(S')\times\mathcal B(k[\operatorname{Ker}f]).
\end{tikzcd}
\]
所以 $\mathcal A(S')$、$\mathcal B(S')$ 分别是 $\mathcal A(S)$、$\mathcal B(S)$ 上的\Term{principal-homogeneous-set}{主齐性集}，\Term{structure-group}{结构群}为 $\mathcal A(k[\operatorname{Ker}f])\xrightarrow\sim\mathcal B(k[\operatorname{Ker}f])$。因此若 $\varphi(S)$ 为\Term{bijection}{双射}，则 $\varphi(S')$ 为\Term{injectivity}{单射}，故为\Term{bijection}{双射}。
\end{proof}
\begin{SGAresult}{推论}{2.16}
设 $\mathcal A$ 是 $\mathcal C_\Lambda$ 上的\Term{homogeneous}{齐性群胚}，且 $\operatorname{Aut}_{\mathcal A(k[\epsilon])}(1_\epsilon)$ \Term{finite-dimensional}{有限维}。则对任意 $\xi\in\operatorname{ob}(\widehat{\mathcal A})$，\Term{functor}{函子}
\[\operatorname{Isom}_\xi:\mathcal C_{R\widehat\otimes_\Lambda R}\longrightarrow(\mathrm{Ens}),\qquad R=P(\xi),\]
是\Term{prorepresentable}{可预表的}。
\end{SGAresult}
\begin{proof}
假定 $\operatorname{Isom}_\xi$ 为\Term{homogeneous}{齐性的}。由 1.11，有一个\Term{minimally-smooth}{极小光滑}\Term{functor}{函子} $\widetilde\varphi:\underline h_{R_1}\to\operatorname{Isom}_\xi$，由 2.15，它是 $\mathcal C_{R\widehat\otimes_\Lambda R}$-\Term{equivalence}{等价}，即 $\operatorname{Isom}_\xi$ \Term{prorepresentable}{可预表}。所以只须证明 $\operatorname{Isom}_\xi$ 为\Term{homogeneous}{齐性的}。考虑 $\mathcal C_{R\widehat\otimes_\Lambda R}$ 中的图 $T_1\to T\leftarrow T_2$，其中 $T_2\to T$ \Term{surjectivity}{满射}。由于 $\mathcal A$ 为\Term{homogeneous}{齐性的}，$\xi_i(T_1)\times_{\xi_i(T)}\xi_i(T_2)$ 在 $\mathcal A$ 中存在，并且对 $i=1,2$ 有\Term{canonical-isomorphism}{典范同构}
\[\xi_i(T_1\times_TT_2)\simeq\xi_i(T_1)\times_{\xi_i(T)}\xi_i(T_2)\quad\text{（参见 2.6(a)）。}\]
因此 $\operatorname{Isom}_\xi(T_1\times_TT_2)\to\operatorname{Isom}_\xi(T_1)\times_{\operatorname{Isom}_\xi(T)}\operatorname{Isom}_\xi(T_2)$ 为\Term{isomorphism}{同构}。
\end{proof}
\begin{SGAresult}{\Term{theorem}{定理}}{2.17}
设 $\mathcal A$ 是 $\mathcal C_\Lambda$ 上的\Term{cofibered-groupoid}{余纤维群胚}。则它能由一个（\Term{normalized-cogroupoid}{规范化}）\Term{smooth-cogroupoid}{光滑余群胚}\Term{prorepresentable}{预表示}，当且仅当
\SourcePage{77}
(1) $\mathcal A$ 为\Term{homogeneous}{齐性的}；
(2) $[\mathcal A(k[\epsilon])]$ 和 $[\operatorname{Aut}(1_\epsilon)]$ 都是\Term{finite-dimensional}{有限维} $k$-\Term{vector-space}{向量空间}。
\end{SGAresult}
\begin{proof}
假设 $\mathcal A$ 由\Term{smooth-cogroupoid}{光滑余群胚}\Term{prorepresentable}{预表示}，即它与 $\widehat{\mathcal C}_\Lambda$ 中某个\Term{smooth-cogroupoid}{光滑余群胚} $(R_0,R_1,\ldots)$ 对应的 $\underline h=\underline h_{(R_0,R_1,\ldots)}$ 为 $\mathcal C_\Lambda$-\Term{equivalence}{等价}。由 2.12(a)，$\underline h$ 为\Term{homogeneous}{齐性的}，并且
\[
\begin{aligned}
[\underline h(k[\epsilon])]
 &=\operatorname{Coker}\bigl(\operatorname{Hom}_{\Lambda\text{-alg}}(R_1,k[\epsilon])
 \underset{s_2}{\overset{s_1}{\rightrightarrows}}
 \operatorname{Hom}_{\Lambda\text{-alg}}(R_0,k[\epsilon])\bigr),\\
\operatorname{Aut}_{\underline h}(1_\epsilon)
 &=\operatorname{Ker}\bigl(\operatorname{Hom}_{\Lambda\text{-alg}}(R_1,k[\epsilon])
 \xrightarrow{s_1}\operatorname{Hom}_{\Lambda\text{-alg}}(R_0,k[\epsilon])\bigr)\\
 &=\operatorname{Ker}\bigl(\operatorname{Hom}_{\Lambda\text{-alg}}(R_1,k[\epsilon])
 \xrightarrow{s_2}\operatorname{Hom}_{\Lambda\text{-alg}}(R_0,k[\epsilon])\bigr)
\end{aligned}
\]
都是\Term{finite-dimensional}{有限维} $k$-\Term{vector-space}{向量空间}，所以 $\mathcal A$ 也具有这些性质。
反之，假设 $\mathcal A$ 满足 (1)、(2)。断言它由一个\Term{normalized-cogroupoid}{规范化}\Term{smooth-cogroupoid}{光滑余群胚}\Term{prorepresentable}{预表示}。由于 $\mathcal A$ 为\Term{homogeneous}{齐性的}，且 $[\mathcal A(k[\epsilon])]$ 是\Term{finite-dimensional}{有限维} $k$-\Term{vector-space}{向量空间}，$\mathcal A$ 有一个\Term{minimally-versal}{极小万有}\Term{formal-object}{形式对象} $\xi$。令 $P(\xi)=R_0$，考虑\Term{functor}{函子}
\[\operatorname{Isom}_\xi:\mathcal C_{R_0\widehat\otimes_\Lambda R_0}\longrightarrow(\mathrm{Ens}),\qquad
 T\longmapsto\operatorname{Isom}_{\mathcal A(T)}(\xi_1(T),\xi_2(T)),\]
其中 $\xi_i(T)=\xi(f_i)$，$f_i:R_0\to T$ 表示复合映射
\[R_0\underset{s_2}{\overset{s_1}{\rightrightarrows}}R_0\widehat\otimes_\Lambda R_0\longrightarrow T.\]
由 2.16，此\Term{functor}{函子}\Term{prorepresentable}{可预表}，即存在
$R_1\in\operatorname{ob}(\widehat{\mathcal C}_{R_0\widehat\otimes_\Lambda R_0})$ 及\Term{isomorphism}{同构} $\varphi:\xi_1(R_1)\xrightarrow\sim\xi_2(R_1)$，使 $\widetilde\varphi:h_{R_1}\xrightarrow\sim\operatorname{Isom}_\xi$ 为 $\mathcal C_{R_0\widehat\otimes_\Lambda R_0}$ 上的\Term{isomorphism}{同构}。于是 $(R_0,R_1,\ldots)$ 定义 $\widehat{\mathcal C}_\Lambda$ 中的\Term{cogroupoid}{余群胚}，且\Term{functor}{函子}
$(\xi,\varphi):\underline h_{(R_0,R_1,\ldots)}\to\mathcal A$ 为 $\mathcal C_\Lambda$-\Term{equivalence}{等价}。由于 $\xi$ 为 $\mathcal A$ 的\Term{minimally-versal}{极小万有}\Term{formal-object}{形式对象}，按定义，复合\Term{functor}{函子}
\[\underline h_{(R_0,R_0,\ldots)}\xrightarrow\epsilon
\underline h_{(R_0,R_1,\ldots)}\xrightarrow{(\xi,\varphi)}\mathcal A\]
为\Term{minimally-smooth}{极小光滑的}；又因为 $(\xi,\varphi)$ 为 $\mathcal C_\Lambda$-\Term{equivalence}{等价}，$\epsilon$ 也为\Term{minimally-smooth}{极小光滑的}，即 $(R_0,R_1,\ldots)$ 是 $\widehat{\mathcal C}_\Lambda$ 中的\Term{normalized-cogroupoid}{规范化}\Term{smooth-cogroupoid}{光滑余群胚}。
\end{proof}
\SourcePage{78}
\begin{SGAresult}{推论}{2.18}
\Term{functor}{函子} $F:\mathcal C_\Lambda\to(\mathrm{Ens})$ \Term{prorepresentable}{可预表}，当且仅当它为\Term{homogeneous}{齐性的}且 $\dim F(k[\epsilon])<\infty$。
\end{SGAresult}
\begin{SGAresult}{推论}{2.19}
$\widehat{\mathcal C}_\Lambda$ 中每个\Term{smooth-cogroupoid}{光滑余群胚}在\Term{isomorphism}{同构}意义下都能\Term{normalized-cogroupoid}{规范化}。
\end{SGAresult}
给定 $\widehat{\mathcal C}_\Lambda$ 中\Term{cogroupoid}{余群胚} $(R_0,R_1,\ldots)$，令 $\overline R_1=R_1/J$，其中 $J$ 为 $R_1$ 中由 $\{s_1(r)-s_2(r)\mid r\in R_0\}$ 生成的\Term{ideal}{理想}。于是 $(R_0,\overline R_1,\ldots)$ 对诱导的\Term{composition-law}{合成律}定义 $R_0$ 上的一个\Term{formal-group-scheme}{形式群概形}。设 $\mathcal A$ 是 $\mathcal C_\Lambda$ 上的\Term{cofibered-groupoid}{余纤维群胚}，经由 $(\xi,\varphi)$ 由\Term{normalized-cogroupoid}{规范化}\Term{smooth-cogroupoid}{光滑余群胚} $(R_0,R_1,\ldots)$ \Term{prorepresentable}{预表示}。设
\[\mathcal C_k\lhook\joinrel\longrightarrow\mathcal C_{R_0}\lhook\joinrel\longrightarrow\mathcal C_{R_0\widehat\otimes_\Lambda R_0}\]
是对应于\Term{surjectivity}{满射} $R_0\widehat\otimes_\Lambda R_0\to R_0\to k$ 的\Term{embedding}{嵌入}，并令 $\operatorname{Aut}_\xi,\operatorname{Aut}^0$ 分别为 $\operatorname{Isom}_\xi$ 在 $\mathcal C_{R_0},\mathcal C_k$ 上的\Term{restriction}{限制}\Term{functor}{函子}。换言之，对每个 $S\in\operatorname{ob}(\mathcal C_{R_0})$，$\operatorname{Aut}_\xi(S)$ 是 $\xi(S)$ 的\Term{automorphism-group}{自同构群}；对每个 $S\in\operatorname{ob}(\mathcal C_k)$，$\operatorname{Aut}^0(S)$ 是 $1_S$ 的\Term{automorphism-group}{自同构群}，其中 $1_S$ 为 $1\in\mathcal A(k)$ 沿 $k\to S$ 的\Term{direct-image}{直接像}。由于 $R_1$ \Term{prorepresentable}{预表示} $\mathcal C_{R_0\widehat\otimes_\Lambda R_0}$ 上的\Term{functor}{函子} $\operatorname{Isom}_\xi$，故 $\overline R_1$ \Term{prorepresentable}{预表示} $\mathcal C_{R_0}$ 上的 $\operatorname{Aut}_\xi$，$k\widehat\otimes_{R_0}\overline R_1$ \Term{prorepresentable}{预表示} $\mathcal C_k$ 上的 $\operatorname{Aut}^0$。虽然 $\mathcal A$ 由\Term{smooth-cogroupoid}{光滑余群胚} $(R_0,R_1,\ldots)$ \Term{prorepresentable}{预表示}，但 $(R_0,\overline R_1,\ldots)$ 和 $(k,k\widehat\otimes_{R_0}\overline R_1,\ldots)$ \Term{generic}{一般}都不\Term{smooth-cogroupoid}{光滑}。事实上，有：
\begin{SGAresult}{推论}{2.20}
设 $\mathcal A$ 是 $\mathcal C_\Lambda$ 上的\Term{homogeneous}{齐性}\Term{cogroupoid}{余群胚}，经由 $(\xi,\varphi)$ 由\Term{normalized-cogroupoid}{规范化}\Term{smooth-cogroupoid}{光滑余群胚} $(R_0,R_1,\ldots)$ \Term{prorepresentable}{预表示}。
(a) $(R_0,\overline R_1,\ldots)$ 和 $(k,k\widehat\otimes_{R_0}\overline R_1,\ldots)$ 分别\Term{prorepresentable}{预表示}\Term{functor}{函子} $\operatorname{Aut}_\xi$ 和 $\operatorname{Aut}^0$。
(b) 以下陈述\Term{equivalence}{等价}：
(i) \Term{formal-group-scheme}{形式群概形} $\underline h_{(R_0,\overline R_1,\ldots)}$ 在 $R_0$ 上\Term{formally-smooth}{形式光滑}；
(ii) 对 $\mathcal A$ 中每个\Term{surjectivity}{满射}\Term{arrow}{箭头} $a'\to a$，$\operatorname{Aut}_{P(a')}(a')\to\operatorname{Aut}_{P(a)}(a)$ \Term{surjectivity}{满射}；
(iii) $R_0$ \Term{prorepresentable}{预表示} $\mathcal C_\Lambda$ 上的\Term{functor}{函子} $[\mathcal A]$。
\end{SGAresult}
\SourcePage{79}
\begin{proof}
(a) 已经说明。至于 (b)，(i)$\Leftrightarrow$(ii) 是显然的；由于 $\widetilde\xi:\underline h_{R_0}\to[\mathcal A]$ 为\Term{minimally-smooth}{极小光滑的}，(ii)$\Leftrightarrow$(iii) 由 2.7 和 2.15 得到。
\end{proof}
\textbf{注记 2.21。} 设 $\mathcal A$ 是 $\mathcal C_\Lambda$ 上的\Term{cofibered-groupoid}{余纤维群胚}，$\mathcal A|\mathcal C_k$ 是 $p:\mathcal A\to\mathcal C_\Lambda$ 沿 $\mathcal C_k\hookrightarrow\mathcal C_\Lambda$ 的拉回。容易看出，若 $\mathcal A$ 经由 $(\xi,\varphi)$ 由 $(R_0,R_1,\ldots)$ \Term{prorepresentable}{预表示}，则 $\mathcal A|\mathcal C_k$ 经由 $(k\widehat\otimes_\Lambda\xi,k\widehat\otimes_\Lambda\varphi)$ 由 $(k\widehat\otimes_\Lambda R_0,k\widehat\otimes_\Lambda R_1,\ldots)$ \Term{prorepresentable}{预表示}。这里 $k\widehat\otimes_\Lambda\xi$ 是 $\xi$ 沿映射 $R_0\to k\widehat\otimes_\Lambda R_0$ 的\Term{direct-image}{直接像}，$k\widehat\otimes_\Lambda\varphi$ 是 $\varphi$ 在映射
\[\operatorname{Isom}_{\mathcal A(R_1)}(\xi(s_1),\xi(s_2))\longrightarrow
\operatorname{Isom}_{\mathcal A(k\widehat\otimes_\Lambda R_1)}
(k\widehat\otimes_\Lambda\xi(s_1),k\widehat\otimes_\Lambda\xi(s_2))\]
下的像。若 $(R_0,R_1,\ldots)$ 为\Term{normalized-cogroupoid}{规范化的}，则 $(k\widehat\otimes_\Lambda R_0,k\widehat\otimes_\Lambda R_1,\ldots)$ 也是如此。$R_1$ 在 $R_0$ 上的\Term{relative-dimension}{相对维数}等于 $k\widehat\otimes_\Lambda R_1$ 在 $k\widehat\otimes_\Lambda R_0$ 上的\Term{relative-dimension}{相对维数}，因为 $R_1$ 在 $R_0$ 上\Term{smooth}{光滑}。

\textbf{注记 2.22。} 设 $\mathcal A$ 由\Term{smooth-cogroupoid}{光滑余群胚} $(R_0,R_1,\ldots)$ 和 $(R_0',R_1',\ldots)$ \Term{prorepresentable}{预表示}，并假设 $(R_0,R_1,\ldots)$ 为\Term{normalized-cogroupoid}{规范化的}。于是得到\Term{homomorphism}{同态} $(i_0,i_1):(R_0,R_1,\ldots)\to(R_0',R_1',\ldots)$。由 2.10(b)，$R_0'$ 是 $R_0$ 上的\Term{formal-power-series-ring}{形式幂级数环}，\Term{homomorphism}{同态} $(i_0,i_1)$ \Term{induction-isomorphism}{诱导同构}
\[R_0'\widehat\otimes_{R_0}(R_0,R_1,\ldots)\xrightarrow\sim(R_0',R_1',\ldots).\]
